\documentclass[11pt]{article}
\usepackage{preamble}
\setlength{\parskip}{4pt}

\begin{document}

\daybanner{AP Statistics \textperiodcentered\ Unit 2 \textperiodcentered\ Topic 2.2 \textperiodcentered\ B: 2026-10-16 \textperiodcentered\ E: 2026-10-30 \textperiodcentered\ TEACHER KEY}{Which Pass Rate Answers Which Question?}

\sectionbanner{CED TETHER}

{\small
\begin{itemize}[leftmargin=1.5em,itemsep=2pt,topsep=2pt]
  \item \textbf{Skill 2.C} --- Calculate summary statistics, relative positions of points within a distribution, correlation, and predicted response.
  \item \textbf{Skill 2.D} --- Compare distributions or relative positions of points within a distribution.
  \item \textbf{Enduring Understanding UNC-1} --- Graphical representations and statistics allow us to identify and represent key features of data.
  \item \textbf{LO UNC-1.Q} --- Calculate statistics for two categorical variables. [Skill 2.C]
  \item \textbf{EK UNC-1.Q.1} --- The marginal relative frequencies are the row and column totals in a two-way table divided by the total for the entire table.
  \item \textbf{EK UNC-1.Q.2} --- A conditional relative frequency is a relative frequency for a specific part of the contingency table (e.g., cell frequencies in a row divided by the total for that row).
  \item \textbf{LO UNC-1.R} --- Compare statistics for two categorical variables. [Skill 2.D]
  \item \textbf{EK UNC-1.R.1} --- Summary statistics for two categorical variables can be used to compare distributions and/or determine if variables are associated.
\end{itemize}
}
\textbf{Essential question:} How do the denominator and study design limit what a percentage can say?\par
\textbf{DOK-3 skill:} 4.B \textperiodcentered\ \textbf{FRQ pattern:} {\small\texttt{marginal-\allowbreak{}vs-\allowbreak{}conditional-\allowbreak{}association-\allowbreak{}is-\allowbreak{}not-\allowbreak{}cause}}\par
\textbf{DOK rationale:} Part (c) weighs competing claims using the day's numerical or graphical evidence and explains a limitation or consequence; the judgment requires linked reasoning beyond a calculation.\par

\frameworkphaseheader{Do Now (first take) $\rightarrow$ Explore (video + follow-along) $\rightarrow$ Exit (finish + turn in)}{1 $\rightarrow$ 3}{45}{%
\begin{itemize}[leftmargin=1.5em,itemsep=2pt,topsep=2pt]
  \item Show the board at the bell and collect a one-sentence first take without confirming a claim.
  \item Run the named video follow-along, then allow about 10 minutes for parts (a)--(c).
  \item Ask for evidence linked to a claim. Collect the sheet and hand-score part (c) with E/P/I.
\end{itemize}
}{%
\begin{itemize}[leftmargin=1.5em,itemsep=2pt,topsep=2pt]
  \item Name the group described by each first-take percentage.
  \item Complete the video follow-along about joint, marginal, and conditional relative frequencies.
  \item Finish (a)--(c), revise the first take in the exit line, and turn in the sheet.
\end{itemize}
}{%
\begin{itemize}[leftmargin=1.5em,itemsep=2pt,topsep=2pt]
  \item Which number or visible feature supports your claim?
  \item Why does the other claim go beyond that evidence?
  \item What would you tell someone who chose the other claim?
\end{itemize}
}{Circulate and ask students to connect their choice to evidence. Score part (c) using the three required reasoning elements; do not require dropped extension work.}

\begin{callout}{\IconWarn\ WATCH FOR}{calloutred}
\begin{itemize}[leftmargin=1.5em,itemsep=2pt,topsep=2pt]
  \item A choice with copied numbers but no explanation of how they support it.
  \item Treating a model or average as a guaranteed individual outcome.
\end{itemize}
\end{callout}

\sectionbanner{SCORING PART (c) --- E / P / I}

\textbf{Expected elements}
\begin{itemize}[leftmargin=1.5em,itemsep=2pt,topsep=2pt]
  \item \textbf{reasoning-1} (required) --- Maps the two rates to all students and tutoring users
  \item \textbf{reasoning-2} (required) --- Corrects the mixed-denominator comparison with 88 versus 70.7 percent and a 17.3-point group gap
  \item \textbf{reasoning-3} (required) --- Explains why self-selection prevents the causal claim
\end{itemize}

{\small\renewcommand{\arraystretch}{1.25}
\noindent\begin{tabular}{|p{0.5in}|p{5.9in}|}
\hline
\textbf{E} & Includes all three required reasoning elements above, with correct contextual evidence. \\ \hline
\textbf{P} & Includes at least one substantive correct reasoning link above, but omits or misstates another required element. \\ \hline
\textbf{I} & Includes no substantive correct reasoning link above; an unsupported choice or copied number alone is insufficient. \\ \hline
\end{tabular}}

\textbf{Common mistakes}
\begin{itemize}[leftmargin=1.5em,itemsep=2pt,topsep=2pt]
  \item Treating 75 percent as the pass rate among nonusers
  \item Subtracting a conditional rate and marginal rate as though they were the two comparison groups
  \item Calling the 17.3-point association causal because the users performed better
\end{itemize}

\clearpage
\focusbanner{TODAY'S PROBLEM --- annotated}

Suppose a school records exam outcomes for 300 students. Students chose for themselves whether to use the tutoring center.\par\smallskip \textbf{Mara:} \emph{The pass rate is 75\%.}\par \textbf{Noah:} \emph{The pass rate is 88\%.}\par \textbf{Elena:} \emph{Tutoring raises the pass rate by 13 percentage points because $88\%-75\%=13$ percentage points.}\par
\begin{center}
\textbf{Exam outcomes for 300 students}\par\smallskip
\begin{tabular}{|l|c|c|c|}
\hline
\textbf{Center?} & \textbf{Passed} & \textbf{Did not pass} & \textbf{Total} \\ \hline
\textbf{Used} & 66 & 9 & 75 \\ \hline
\textbf{Did not} & 159 & 66 & 225 \\ \hline
\textbf{Total} & 225 & 75 & 300 \\ \hline
\end{tabular}
\end{center}
\begin{firsttakebox}{FIRST TAKE (5 min)}
Can 75\% and 88\% both be correct pass rates? State what group you think each percentage describes.\par
\answer{Not graded. Expect some students to believe two different pass rates cannot both be correct until their denominators are named.}
\end{firsttakebox}

\textit{Rules callout ``NAME THE DENOMINATOR, THEN THE CLAIM (keep on the page)'' is printed on the student sheet.}\par\medskip

\dokbadge{1}~\textbf{(a)} Label 75 percent and 88 percent as marginal or conditional. Name the group for each.\par
\answer{The $75\%$ is the marginal relative frequency of passing among all 300 students. The $88\%$ is the conditional relative frequency of passing among students who used the tutoring center; the condition is used the center.}

\dokbadge{2}~\textbf{(b)} Verify the 75 percent and 88 percent rates. The nonuser pass rate is about 70.7 percent; find the user-minus-nonuser gap.\par
\answer{Overall: $225/300=75\%$. Users: $66/75=88\%$. The group gap is about $88-70.7=17.3$ percentage points.}

\dokbadge{3}~\textbf{(c)} Can Mara and Noah both be right? Explain their denominators. Use the two group rates to correct Elena's 13-point comparison, and explain why her word raises goes beyond these data. Write 3--4 sentences.\par
\answer{Both rates are correct: Mara uses all 300 students, while Noah uses the 75 tutoring users. Users passed at $88\%$ versus about $70.7\%$ for nonusers, so the observed group gap is about 17.3 percentage points. Elena subtracts the overall rate from the user rate; the overall group includes users and nonusers, so 13 points is not the direct group gap. Students chose tutoring, so differences such as preparation could help explain the association; the table does not show tutoring raised pass rates.}

\begin{summaryexitbox}{EXIT}
One thing the video changed about my first take: \hrulefill
\end{summaryexitbox}

\end{document}
